\section{Thiết lập thực nghiệm}
\label{sec:setup}

Mục này mô tả cách tái lập SyscallAD để \emph{chạy lại} thực nghiệm. Số liệu kết quả
được báo cáo ở Mục~\ref{sec:results} \emph{chỉ} lấy từ lần chạy lại trên bộ dữ liệu
DongTing đầy đủ (chưa phải số liệu trong mục này).

\subsection{Nguyên tắc: mọi công cụ chạy bằng Docker}
Để bảo đảm khả năng tái lập và không làm ``bẩn'' máy chủ, toàn bộ chuỗi công cụ
(huấn luyện, đánh giá, biên dịch báo cáo) đều chạy trong container. Quy trình huấn
luyện và đánh giá được đóng gói bằng \texttt{docker-compose} và điều phối bởi
\texttt{Experiment/run\_experiment.sh}.

\subsection{Huấn luyện ngoại tuyến SyscallAD}
Một nguồn chân lý duy nhất là \texttt{train.py} (mọi hàm học máy, không vẽ đồ thị);
sổ tay \texttt{train\_dongting.ipynb} nhập lại từ \texttt{train.py} và bổ sung đồ thị
+ bình luận. Hai chế độ luôn đồng bộ.

\begin{lstlisting}[language=bash]
cd DeSFAM/training
# Chay nhanh, khong tuong tac:
docker compose -f docker-compose.pipeline.yml run --rm train
# Hoac tuong tac (Jupyter Lab :8888):
docker compose -f docker-compose.pipeline.yml --profile interactive up jupyter
\end{lstlisting}

Đầu ra (lưu vào \texttt{../model/}): bộ chuẩn hóa đặc trưng, mô hình iForest, trọng
số encoder/decoder VAE, các bộ chuẩn hóa điểm thành phần, tham số tập hợp
($\alpha=0{,}7$), từ vựng đặc trưng (\texttt{fe\_report.json}) và các chỉ số.

\subsection{Cấu hình thí nghiệm (khai báo trước)}
Bảng~\ref{tab:config} là các \emph{siêu tham số đã chọn} (cố định trước khi chạy);
các đại lượng phụ thuộc dữ liệu (số chuỗi, số cửa sổ mỗi tập) sẽ được điền từ lần
chạy lại trên DongTing đầy đủ.

\begin{table}[htbp]
  \centering
  \caption{Cấu hình thí nghiệm đã khai báo cho SyscallAD.}
  \label{tab:config}
  \begin{tabular}{ll}
    \toprule
    \textbf{Tham số} & \textbf{Giá trị} \\
    \midrule
    Thí nghiệm & \texttt{dongting\_15\_3} \\
    Cửa sổ trượt (độ dài, bước) & $(15,\ 3)$ \\
    Số chiều đặc trưng & 43 \\
    VAE \texttt{latent\_dim} / \texttt{hidden\_dim} & 8 / 32 \\
    iForest \texttt{contamination} & 0{,}02 \\
    Trọng số hợp nhất $\alpha$ & 0{,}7 \\
    Bảng chữ cái syscall an ninh & 23 syscall an ninh (\texttt{SENSOR\_ALPHABET=1}) \\
    Đặc trưng thời gian & không (\texttt{has\_temporal=false}) \\
    Chiến lược ngưỡng & \texttt{p995} (phân vị 99,5 lỗi kiểm định lành tính) \\
    \midrule
    Số chuỗi DongTing (nạp được / công bố) & 18\,874 / 18\,966 \\
    Số cửa sổ train (lành tính) / val / test & 90\,456 / 896\,581 / 687\,926 \\
    Số mẫu lành tính (PrefixSpan) & 150 \\
    \bottomrule
  \end{tabular}
\end{table}

\subsection{Giao thức đánh giá và tiêu chí tái lập thành công}
\label{subsec:protocol}
Để tránh diễn giải tùy tiện sau khi có số liệu, báo cáo \emph{cố định trước} giao
thức đánh giá:
\begin{itemize}[leftmargin=1.6em]
  \item \textbf{Đơn vị đánh giá:} mỗi \emph{cửa sổ trượt} $(15,3)$ là một mẫu. Nhãn
        cửa sổ kế thừa nhãn chuỗi DongTing chứa nó. \emph{Cảnh báo phương pháp:} một
        chuỗi tấn công chứa nhiều cửa sổ ``trông lành tính'' (không có syscall đặc
        quyền nào) nhưng vẫn bị gán nhãn tấn công; điều này gây nhiễu nhãn và có thể
        thổi phồng/làm lệch AUC/AP/F1. Do đó báo cáo \emph{đồng thời} báo cáo (a) chỉ
        số mức \emph{cửa sổ} và (b) chỉ số mức \emph{chuỗi} (một chuỗi bị gắn cờ nếu
        có cửa sổ vượt ngưỡng), và nêu tỉ lệ cửa sổ tấn công thực sự chứa syscall đặc
        quyền.
  \item \textbf{Phân chia dữ liệu:} dùng đúng phân chia
        \texttt{DTDS-train}/\texttt{validation}/\texttt{test} của DongTing; huấn
        luyện trên cửa sổ \emph{bình thường} của tập train, chọn ngưỡng trên tập
        validation, báo cáo số trên tập test. Nêu rõ rằng việc chọn hạt giống VAE
        theo AUC kiểm định \emph{có dùng nhãn tấn công của tập validation}; tập test
        tuyệt đối không tham gia chọn mô hình/ngưỡng.
  \item \textbf{Chỉ số:} AUC và Average Precision (AP)---đo khả năng \emph{xếp hạng}
        bất thường, độc lập ngưỡng; F1/Precision/Recall---đo \emph{điểm chốt nhị
        phân} tại ngưỡng đã chọn. Báo cáo kèm \emph{trung bình $\pm$ độ lệch chuẩn}
        qua các hạt giống thay vì một điểm số đơn lẻ.
  \item \textbf{Tiêu chí ``tái lập đạt'' (không neo theo số liệu công bố theo mô
        hình---vì DeSFAM \emph{không} công bố bảng AUC/AP/F1 theo mô hình trên
        DongTing):} coi là đạt nếu (a) bộ trích đặc trưng dùng khi huấn luyện và khi
        suy luận sinh véc-tơ \emph{trùng khít} (kiểm thử ngang hàng đạt); (b) khả năng \emph{xếp
        hạng} cao và ổn định (AUC/AP rõ rệt trên mức ngẫu nhiên, ổn định qua hạt
        giống); (c) định hướng \emph{precision cao} (ít báo động giả) được giữ.
        F1/recall phụ thuộc ngưỡng nên \emph{không} dùng làm tiêu chí đạt/không; mọi
        đối chiếu với DeSFAM chỉ ở mức \emph{định tính} (Mục~\ref{sec:results}).
  \item \textbf{Thí nghiệm phân định (bắt buộc):} để tách ``phát hiện bất thường theo
        chuỗi'' khỏi ``chỉ phát hiện sự hiện diện của syscall đặc quyền'', báo cáo
        chạy một \emph{ablation}: so AUC của mô hình đầy đủ với một bộ phân loại
        \emph{chỉ dùng đặc trưng \texttt{disc}} (Mục~\ref{subsec:confound}). Nếu hai
        bên gần nhau, kết luận phải lùi về ``phát hiện theo hiện diện là tái lập
        được; mô hình học máy bổ sung biên [đo được]''.
\end{itemize}

\paragraph{Đối sánh mã syscall giữa hai lớp.} Trong DongTing, chuỗi lành tính lưu
dưới dạng \emph{ID số}, còn chuỗi tấn công lưu dưới dạng \emph{tên} syscall. Báo cáo
dùng \texttt{syscall\_64.tbl} (bỏ dòng ABI x32) để ánh xạ tên\,$\leftrightarrow$\,ID
nhất quán và kiểm tra rằng cả 23 syscall trong bảng chữ cái khôi phục \emph{giống
hệt} từ cả hai cách lưu---nếu lệch, các bit \texttt{disc} sẽ bật bất đối xứng giữa
hai lớp và tạo/triệt tiêu độ phân tách một cách giả tạo.

\paragraph{Trung thành với nguồn so với thích nghi.} Bảng~\ref{tab:fidelity} tách
bạch phần tái lập \emph{đúng} công bố với phần \emph{thích nghi có lý do}, giúp diễn
giải đúng chênh lệch số liệu ở Mục~\ref{sec:results}.

\begin{table}[htbp]
  \centering
  \caption{Trung thành với DeSFAM so với thích nghi trong bản tái lập.}
  \label{tab:fidelity}
  \small
  \begin{tabularx}{\linewidth}{@{}>{\raggedright\arraybackslash}p{4.6cm} l >{\raggedright\arraybackslash}X@{}}
    \toprule
    \textbf{Hạng mục} & \textbf{Trạng thái} & \textbf{Ghi chú} \\
    \midrule
    VAE+iForest, $\alpha=0{,}7$, cửa sổ $(15,3)$ & tái lập đúng & $\alpha$ \emph{kế thừa}, không tinh chỉnh lại \\
    Huấn luyện một lớp (chỉ lành tính)           & tái lập đúng & \\
    Đặc trưng thời gian ($\Delta t$)             & thích nghi (tắt) & DongTing không có dấu thời gian \\
    Bảng chữ cái syscall                         & thích nghi (23 lời gọi) & tập syscall an ninh liên quan leo thang \\
    Ngưỡng                                       & thích nghi (cố định) & dùng thay EMA để dễ đối chiếu \\
    Mô-đun 1, 3                                  & ngoài phạm vi & không hiện thực hóa lại \\
    \bottomrule
  \end{tabularx}
\end{table}

\subsection{Dữ liệu tấn công và lành tính}
Báo cáo đánh giá hoàn toàn \emph{ngoại tuyến} trên DongTing: chuỗi \emph{lành tính}
đến từ bốn bộ kiểm thử hồi quy nhân Linux, chuỗi \emph{tấn công} từ các chương trình
kích hoạt lỗi (syzkaller) trên 26 phiên bản nhân---các chuỗi này phản ánh các họ
khai thác leo thang đặc quyền/thoát container (ví dụ CVE-2021-4034, CVE-2022-0847,
CVE-2019-5736). Không có thành phần phục vụ trực tiếp; mọi số liệu lấy từ tập kiểm
tra DongTing.
